\section{Cơ sở lý thuyết}

\subsection{Trí tuệ nhân tạo tạo sinh và mô hình ngôn ngữ lớn trong giáo dục}

\subsubsection{Khái niệm và nguyên lý hoạt động của trí tuệ nhân tạo tạo sinh}
Trí tuệ nhân tạo tạo sinh là một lĩnh vực của trí tuệ nhân tạo tập trung vào việc tạo lập các nội dung mới như văn bản, hình ảnh hoặc mã nguồn dựa trên dữ liệu huấn luyện có sẵn. Trong xử lý ngôn ngữ tự nhiên, sự phát triển của các mô hình ngôn ngữ lớn đã thay đổi phương thức tương tác giữa con người và máy tính. Bản chất của các mô hình ngôn ngữ lớn là mạng nơ-ron học sâu được huấn luyện trên các tập dữ liệu văn bản có quy mô lớn. Nguyên lý hoạt động cốt lõi của các mô hình này là dự đoán phần tử ngôn ngữ tiếp theo trong một chuỗi văn bản dựa trên xác suất xuất hiện của ngữ cảnh đứng trước. Quá trình xử lý văn bản được thực hiện bằng cách phân rã câu chữ thành các đơn vị thẻ từ. Khả năng duy trì thông tin trong một cuộc hội thoại phụ thuộc vào độ rộng của cửa sổ ngữ cảnh, tức là số lượng thẻ từ tối đa mà mô hình có thể tiếp nhận và xử lý tại một thời điểm. Cửa sổ ngữ cảnh càng lớn giúp mô hình tổng hợp thông tin từ các tài liệu dài một cách logic và nhất quán hơn.

\subsubsection{Prompt engineering trong giáo dục đại học}
\textit{Prompt engineering} là kỹ thuật thiết kế và tinh chỉnh các câu lệnh đầu vào nhằm định hình phản hồi của mô hình ngôn ngữ lớn theo mục tiêu cụ thể. Trong môi trường giáo dục đại học, \textit{prompt engineering} đóng vai trò quan trọng quyết định chất lượng học liệu do công nghệ tạo ra. Việc thiết kế câu lệnh bao gồm việc phân định rõ câu lệnh hệ thống để xác định vai trò chuyên môn của mô hình, các kỹ thuật gợi ý kèm ví dụ mẫu để định hình định dạng dữ liệu đầu ra, và kỹ thuật đóng vai nhằm mô phỏng tư duy của một chuyên gia sư phạm. Áp dụng đúng các kỹ thuật này giúp giảng viên kiểm soát cấu trúc thông tin đầu ra, tránh các phản hồi chung chung và tập trung vào các kiến thức kỹ thuật chuyên sâu.

\subsubsection{Ứng dụng thực tiễn của trí tuệ nhân tạo tạo sinh trong soạn thảo học liệu}
Trí tuệ nhân tạo tạo sinh hỗ trợ giảng viên tự động hóa nhiều công đoạn chuẩn bị tài liệu giảng dạy. Đối với khâu xử lý tài liệu học thuật, công nghệ có khả năng tóm tắt các giáo trình hoặc bài báo khoa học dài thành các luận điểm cốt lõi. Trong thiết kế bài giảng, mô hình hỗ trợ xây dựng kịch bản giảng dạy, phân rã nội dung bài giảng thành cấu trúc dạng danh sách để đưa lên bản trình chiếu và gợi ý các tình huống thực tiễn. Ở khâu kiểm tra đánh giá, các mô hình ngôn ngữ lớn hỗ trợ biên soạn ngân hàng câu hỏi trắc nghiệm hoặc tự luận từ tài liệu đầu vào một cách nhanh chóng. Việc tự động hóa các tác vụ hành chính này giúp giảng viên tiết kiệm thời gian chuẩn bị và tập trung vào các hoạt động tương tác với người học.

\subsubsection{Những giới hạn và thách thức học thuật}
Bên cạnh các lợi ích, việc ứng dụng trí tuệ nhân tạo tạo sinh trong môi trường học thuật đối mặt với nhiều giới hạn. Thách thức kỹ thuật lớn nhất là hiện tượng ảo tưởng thông tin, khi mô hình đưa ra câu trả lời sai lệch nhưng có văn phong thuyết phục. Trong đào tạo khối ngành kỹ thuật, sai số này có thể ảnh hưởng đến tính chính xác của bài giảng. Ngoài ra, việc tải tài liệu giảng dạy lên hệ thống trực tuyến cũng đặt ra rủi ro về an toàn thông tin và bản quyền học thuật. Cuối cùng, việc xử lý thuật ngữ chuyên ngành tiếng Việt của một số mô hình ngôn ngữ lớn hiện nay còn hạn chế, đòi hỏi giảng viên phải thực hiện quy trình kiểm duyệt nghiêm ngặt đối với sản phẩm đầu ra để bảo đảm tính liêm chính học thuật.

\subsection{Thang đo nhận thức Bloom và đánh giá kết quả học tập}

\subsubsection{Cấu trúc thang đo nhận thức Bloom cải tiến}
Thang đo nhận thức Bloom cải tiến là khung lý thuyết phân loại các mục tiêu học tập và mức độ tư duy thành sáu cấp độ nhận thức từ thấp đến cao. Cấp độ thứ nhất là nhớ, liên quan đến khả năng nhận biết và tái hiện các định nghĩa hoặc dữ kiện đã học. Cấp độ thứ hai là hiểu, yêu cầu người học giải thích ý nghĩa, phân loại hoặc tóm tắt thông tin. Cấp độ thứ ba là áp dụng, đòi hỏi việc vận dụng kiến thức lý thuyết để giải quyết các bài tập hoặc tình huống cụ thể. Cấp độ thứ tư là phân tích, yêu cầu chia nhỏ thông tin để làm rõ cấu trúc và mối quan hệ giữa các thành phần. Cấp độ thứ sáu là đánh giá, đòi hỏi người học đưa ra nhận định hoặc phê phán giải pháp dựa trên các tiêu chí chuẩn hóa. Cấp độ cao nhất là sáng tạo, yêu cầu tổng hợp các yếu tố để xây dựng phương án hoặc sản phẩm mới.

\subsubsection{Vai trò của thang Bloom trong thiết kế học liệu đánh giá}
Trong đào tạo đại học khối ngành kỹ thuật, thang đo Bloom là cơ sở để thiết kế công cụ kiểm tra đánh giá chuẩn xác năng lực học tập của sinh viên. Học liệu đánh giá cần phân bổ đều các câu hỏi từ bậc thấp đến bậc cao, thay vì chỉ tập trung vào kiểm tra khả năng nhớ máy móc của sinh viên. Việc biên soạn các câu hỏi ở bậc áp dụng, phân tích hoặc đánh giá giúp kiểm tra tư duy kỹ thuật của sinh viên trước các bài toán thực tế như phân tích nguyên nhân sự cố hệ thống hay lựa chọn giải pháp công nghệ tối ưu.

\subsubsection{Khả năng tích hợp trí tuệ nhân tạo để sinh câu hỏi theo thang Bloom}
Việc ứng dụng trí tuệ nhân tạo tạo sinh cho phép tự động hóa quy trình xây dựng câu hỏi đánh giá bám sát thang đo Bloom. Thông qua các câu lệnh gợi ý có cấu trúc, giảng viên có thể yêu cầu mô hình sử dụng các động từ hành động đặc trưng cho từng cấp độ nhận thức để biên soạn câu hỏi. Cách tiếp cận này giúp rút ngắn quy trình xây dựng ngân hàng câu hỏi, đồng thời bảo đảm các câu hỏi kiểm tra tuân thủ đúng mục tiêu sư phạm của chương trình đào tạo.

\subsection{Thuyết kiến tạo và thiết kế hoạt động dạy học chủ động}

\subsubsection{Nguyên lý cơ bản của thuyết kiến tạo}
Thuyết kiến tạo là một quan điểm học tập khẳng định người học chủ động xây dựng tri thức cho bản thân thông qua trải nghiệm thực tế và quá trình tương tác với môi trường xung quanh, thay vì chỉ tiếp thu thụ động từ giảng viên. Trong không gian lớp học kiến tạo, vai trò của giảng viên dịch chuyển từ người truyền đạt kiến thức trực tiếp sang người định hướng và tổ chức các hoạt động học tập. Học tập được nhìn nhận như một quá trình liên tục điều chỉnh các mô hình nhận thức của người học dựa trên những trải nghiệm mới.

\subsubsection{Thiết kế hoạt động dạy học tích cực}
Áp dụng thuyết kiến tạo đòi hỏi giảng viên thiết kế các hoạt động học tập tích cực để kích thích sự tham gia của sinh viên. Các phương pháp phổ biến bao gồm học tập qua giải quyết vấn đề, học tập qua dự án, nghiên cứu tình huống thực tế và mô hình lớp học đảo ngược. Trong các mô hình này, sinh viên được yêu cầu tự tìm hiểu kiến thức nền tảng trước giờ học, để thời gian trên lớp được dành cho các hoạt động thảo luận nhóm, giải bài tập tình huống hoặc thực hành ứng dụng dưới sự hướng dẫn của giảng viên.

\subsubsection{Ứng dụng trí tuệ nhân tạo tạo sinh trong đề xuất hoạt động dạy học kiến tạo}
Trí tuệ nhân tạo tạo sinh hỗ trợ giảng viên thiết kế kịch bản hoạt động dạy học dựa trên các nguyên lý của thuyết kiến tạo. Mô hình có thể gợi ý các tình huống thực tiễn gắn liền với nội dung bài học, đề xuất các câu hỏi thảo luận mở nhằm kích thích phản biện, hoặc xây dựng các bài tập nhóm giải quyết vấn đề. Những kịch bản hoạt động do công nghệ gợi ý giúp bài giảng phong phú hơn, khuyến khích sinh viên tự khám phá kiến thức và phát triển tư duy thực hành kỹ thuật.
